\adjustbox{max width=\textwidth}{%
\begin{tabular}{lccccc}
\hline  & Treatment & Control & Normalized & \it T-test & \\
 & \it Mean & \it Mean & Difference & (\it p-value) & \\
 & (1) & (2) & (3) & (4) & \\
\hline 
N. of subdistricts & 62 & 57 &  &  & \\
Population (millions, UN, 2020 estimate) & 0.737 & 0.654 & 0.103 & 0.427 & \\
Share of rural population (UN, 2020 estimate) & 0.671 & 0.712 & -0.125 & 0.336 & \\
Area (thousand sq km) & 0.060 & 0.067 & -0.128 & 0.320 & \\
Elevation mean (m) & 201.598 & 204.015 & -0.016 & 0.905 & \\
Elevation std. deviation (m) & 11.480 & 13.916 & -0.095 & 0.462 & \\
Distance to district capital (thousand km) & 2.021 & 1.731 & 0.206 & 0.112 & \\
Malaria tests conducted per thousand people & 2.489 & 3.754 & -0.131 & 0.312 & \\
Urban malaria positivity rate (\%) & 2.175 & 0.978 & 0.165 & 0.385 & \\
Rural malaria positivity rate (\%) & 2.666 & 0.621 & 0.153 & 0.273 & \\
\hline \\[-5.1ex] &  &  &  &  & \\
\hline
\multicolumn{5}{p{.9\linewidth}}{\scriptsize Notes: The table reports subdistrict characteristics before the campaign, by experimental condition. Only subdistricts with more than 50\% of the population within study circles are included. The normalized difference is the difference between the two means divided by the square root of the sum of the variances of the two variables.}
\end{tabular}
}%
\\